\documentclass[10pt,a4paper]{amsart}
\usepackage[margin=19mm]{geometry}
\usepackage{amsmath,amssymb,mathtools}
\usepackage[scr=rsfs,cal=euler]{mathalfa}
\usepackage{xfrac}
\usepackage[unicode,hidelinks,bookmarks=false]{hyperref}
\newcommand{\ie}{i.e.\ }
\newcommand{\cf}{cf.\ }
\newcommand{\resp}{respectively\ }
\newcommand{\Subsection}{\subsection}
\providecommand{\nobreakdash}{\nobreak\hbox{-}}
\setlength{\emergencystretch}{2em}
\multlinegap0pt
\usepackage{latexsym}
\allowdisplaybreaks
\usepackage{bbm}
\usepackage{mathtools}
\usepackage{ulem}
\usepackage{tikz-cd}
\usepackage{mathrsfs}
\usepackage{todonotes}
\usepackage{enumitem}
\usepackage{verbatim}
\usepackage{amsaddr}
\usepackage{dsfont}
\newtheorem{thm}{Theorem}[section]
\newtheorem*{thm*}{Theorem}
\newtheorem{Con}[thm]{Conjecture}
\newtheorem{cor}{Corollary}
\newtheorem{prop}{Lemma}
\newtheorem{lem}[thm]{Lemma}
\newtheorem{pro}[thm]{Lemma}
\newtheorem{defin}{Definition}
\newtheorem{rem}{Remark}[section]
\newtheorem{exa}[thm]{Example}
\newtheorem*{xrem}{Remark}
\numberwithin{equation}{section}
\newcommand{\bg}{\big}
\newcommand{\bgl}{\big(}
\newcommand{\bgr}{\big)}
\newcommand{\bgg}{\bigg}
\newcommand{\bggl}{\bigg(}
\newcommand{\bggr}{\bigg)}
\newcommand{\Bg}{\Big}
\newcommand{\Bgl}{\Big(}
\newcommand{\Bgr}{\Big)}
\newcommand{\Bgg}{\Bigg}
\newcommand{\Bggl}{\Bigg(}
\newcommand{\Bggr}{\Bigg)}
\newcommand{\lo}{\log_2}
\newcommand{\lt}{\log_3}
\newcommand{\inv}{^{-1}}
\newcommand{\ph}{\Phi}
\newcommand{\mbc}{\mathbb{C}}
\newcommand{\mbd}{\mathbb{D}}
\newcommand{\mbe}{\mathbb{E}}
\newcommand{\mbt}{\mathbb{T}}
\newcommand{\mbz}{\mathbb{Z}}
\newcommand{\mbr}{\mathbb{R}}
\newcommand{\mbs}{\mathbb{S}}
\newcommand{\mbf}{\mathbb{F}}
\newcommand{\mbh}{\mathbb{H}}
\newcommand{\mbl}{\mathbb{L}}
\newcommand{\mbn}{\mathbb{N}}
\newcommand{\mbp}{\mathbb{P}}
\newcommand{\mbq}{\mathbb{Q}}
\newcommand{\mbb}{\mathbb{B}}
\newcommand{\mbm}{\mathbb{M}}
\newcommand{\mbbx}{\mathbbm{x}}
\newcommand{\mbby}{\mathbbm{y}}
\newcommand{\mch}{\mathcal{H}}
\newcommand{\mca}{\mathcal{A}}
\newcommand{\mcb}{\mathcal{B}}
\newcommand{\mcc}{\mathcal{C}}
\newcommand{\mcd}{\mathcal{D}}
\newcommand{\mce}{\mathcal{E}}
\newcommand{\mcf}{\mathcal{F}}
\newcommand{\mcg}{\mathcal{G}}
\newcommand{\mcj}{\mathcal{J}}
\newcommand{\mcpg}{\mathcal{g}}
\newcommand{\mcph}{\mathcal{h}}
\newcommand{\mcn}{\mathcal{N}}
\newcommand{\mcp}{\mathcal{P}}
\newcommand{\mck}{\mathcal{K}}
\newcommand{\mcl}{\mathcal{L}}
\newcommand{\mco}{\mathcal{O}}
\newcommand{\mcq}{\mathcal{Q}}
\newcommand{\mcm}{\mathcal{M}}
\newcommand{\mcr}{\mathcal{R}}
\newcommand{\mcs}{\mathcal{S}}
\newcommand{\mfa}{\mathfrak{A}}
\newcommand{\mfpa}{\mathfrak{a}}
\newcommand{\mfb}{\mathfrak{B}}
\newcommand{\mfg}{\mathfrak{G}}
\newcommand{\mfr}{\mathfrak{R}}
\newcommand{\mfpg}{\mathfrak{g}}
\newcommand{\mfph}{\mathfrak{h}}
\newcommand{\mfz}{\mathfrak{z}}
\newcommand{\mse}{\mathscr{E}}
\newcommand{\msk}{\mathscr{K}}
\newcommand{\mskt}{\mathscr{K}_\tau}
\newcommand{\msu}{\mathscr{U}}
\newcommand{\mst}{\mathscr{T}}
\newcommand{\mckt}{\mathcal{K}_\tau}
\newcommand{\mmd}{\mathrm{d}}
\newcommand{\mme}{\mathrm{e}}
\newcommand{\mmi}{\mathrm{i}}
\newcommand{\mit}{\mathrm{i}t}
\newcommand{\re}{\operatorname{Re}}
\newcommand{\im}{\operatorname{Im}}
\newcommand{\me}{\operatorname{meas}}
\newcommand{\whp}{\widehat{\Phi}}
\hypersetup{colorlinks=true,linkcolor=blue,anchorcolor=blue,citecolor=blue}
\usepackage{color}
\DeclareMathOperator\dif{d\!}
\newcommand{\newabstract}[1]{%
	\par\bigskip
	\csname otherlanguage*\endcsname{#1}%
	\csname captions#1\endcsname
	\item[\hskip\labelsep\scshape\abstractname.]
}
\begin{document}
\title[Extreme values of the real part of the Riemann zeta function]{Extreme values of the real part of the Riemann zeta function}
\author{Qiyu Yang, Shengbo Zhao}
\date{}
\maketitle
\noindent\emph{Source and licence.} Extreme values of the real part of the Riemann zeta function. Version arXiv:2607.03654v1 (CC BY 4.0). This material is distributed under the Creative Commons Attribution 4.0 International licence, http://creativecommons.org/licenses/by/4.0/, as recorded in the arXiv OAI licence metadata for this exact version and shown on the abstract page https://arxiv.org/abs/2607.03654v1. Full rights notice: licenses/arxiv-2607.03654-CC-BY-4.0.txt.\par\smallskip
\noindent\textit{Editorial scope.} Counted unit: the original \texttt{lem} environment \texttt{pro2} at source lines 327--334 together with its complete proof at lines 335--337; the delivered document keeps source lines 250--338 so that every referenced label and every citation resolves inside it. Native editable source is the paper's own concatenated TeX; the AMS wrapper is editorial formatting only. No publisher class, upstream script or unrelated artwork is redistributed.
\par\smallskip\noindent\textit{Reference note.} This article ships no compiled bibliography (biblatex); the entries delivered here are composed from the author's own BibTeX (.bib) fields, with no field invented and no entry dropped.
\section{Preliminaries}
\label{secpre}

\subsection{Constructing the resonator}
\label{secConstucting}

In this subsection, we construct the resonator \(R(t)\) to employ the resonance method. Our construction follows that of Bondarenko and Seip \cite{bondarenko2018argument}.
\par
For large \(T\), let \(N = \lfloor T^\kappa \rfloor\) be a large integer, where \(\kappa \in (0,1)\). Let \(\eta \in (0,1)\) be a parameter to be chosen later. Then, we define 
\[
\mcp = \bg\{p : \mme \log N \log_{2}N < p \leq \exp\bg((\log_2 N)^{\eta}\bg) \log N \log_{2}N \bg\}.
\]
Next, we define the multiplicative function \(f(n)\) to be supported on the set of square-free numbers, with values for \(p \in \mcp\) as
\[
f(p)=\sqrt{\frac{\log N \lt N}{\lo N}} \frac{1}{\sqrt{p}(\log p - \lo N-\lt N)}.
\]
Furthermore, if \(p \notin \mcp\), \(f(p)=0\).
\par
For \(k \in \bg\{1,\dots, \lfloor (\log_2 N)^{\eta} \rfloor\bg\}\), define the set \(\mcp_k\) as follows:
\[
\mathcal{P}_{k} = \bg\{p : \mme^k \log N \log_{2}N < p \leq \mme^{k+1} \log N\log_{2}N \bg\}.
\]
Then for fixed \(b \in (1,1/\eta)\), let
\[
\mathcal{M}_{k} = \left\{ n \in \mathrm{supp}(f) : n \text{ has at least } \Delta_k := \frac{b \log N }{k^2 \log_3 N} \text{ prime divisors in } \mathcal{P}_{k} \right\}.
\]
Let \(\mathcal{M}^{\prime}_{k}\) be the set of integers from \(\mathcal{M}_{k}\) that have prime divisors only in \(\mathcal{P}_{k}\), then set
\[
\mathcal{M} = \mathrm{supp}(f) \setminus \bigcup_{k = 1}^{\lfloor(\log_2 N)^{\eta}\rfloor} \mathcal{M}_{k}.
\]
Clearly, \(\mathcal{M}\) is divisor closed. Indeed, if \(m^{\prime} \mid m \in \mathcal{M}\), we have \(m^{\prime} \in \mathcal{M}\). Following the similar argument as in \cite[Lemma 2]{bondarenko2017large}, we have \(|\mcm| \le N\).
\par
Let \(\mch\) be the set of integers \(h\) such that 
\[
\bgg[\Bgl1+\frac{1}{T} \Bgr^h, \Bgl1+\frac{1}{T} \Bgr^{h+1} \bgg] \bigcap \mcm \neq \emptyset.
\]
Then, for all \(h \in \mch\), let \(m_h\) be the minimum of \(\bg[\bg(1+T\inv\bg)^h, \bg(1+T\inv\bg)^{h+1}\bg] \cap \mcm\), and \(\mcm^\prime\) be the set of all \(m_h\). Define \(r(n)\) to satisfy
\[
r(m_h) = \Bigg( \sum_{n \in \mathcal{M},  (1 +T\inv)^{h - 1} \le n \le (1 + T\inv)^{h + 2}} f(n)^2 \Bigg)^{\frac{1}{2}}
\]
for all \(m_h \in \mcm^\prime\). Here, \(f(n)\) denotes the multiplicative function defined previously. Furthermore, we define the resonator
\[
R(t) = \sum_{m \in \mcm^{\prime}}r(m)m^{-\mit}.
\]
Note that we choose \(N = \lfloor T^\kappa \rfloor\) for some \(\kappa \in (0,1)\). Thus, we have
\begin{align}
\label{R0-upper}
    |R(0)|^2 \le |\mcm^\prime|\sum_{m \in \mcm^\prime} r(m)^2 \ll N \sum_{n \in \mcm}f(n)^2.
\end{align}
\par
Finally, as in \cite{bondarenko2017large, bondarenko2018argument}, we take \(\Phi(t) := \mme^{-t^{2}/2}\) to be a Gaussian function. Then \(\ph(t)\) decays rapidly, and its Fourier transform satisfies \(\whp(\xi) = \sqrt{2\pi}\ph(\xi)>0\). Combining the construction of \(\mcm^\prime\) and \(r(n)\) with the Cauchy-Schwarz inequality, we obtain the following upper bound; see \cite[Lemma 5]{bondarenko2018argument}:
\begin{align}
    \label{R2phi-upper}
    \int_{\mbr}|R(t)|^2\Phi \Big( \frac{t}{T} \Big) \mmd t \ll T \sum_{n \in \mcm}f(n)^2.
\end{align}

\subsection{Auxiliary lemmas}
\label{secauxiliary}

In this subsection, we collect two lemmas that play an important role in the subsequent proof. First, we establish a convolution formula for \(\zeta(s)\). The following Lemma can be found in \cite{bondarenko2018argument}.
\begin{lem}[\cite{bondarenko2018argument}, Lemma 1]
    \label{pro1}
    Let \(\sigma \in [1/2,1)\) be fixed. Assume that \(K(x+\mmi y)\) is a holomorphic
function in the strip \(\sigma-2 \le y \le 0\), satisfying the growth condition
\begin{align}
    \label{pro1-growth}
    \max_{\sigma -2\le y \le 0}|K(x+\mmi y)| = O\Bgl \frac{1}{|x|^2+1} \Bgr
\end{align}
when \(|x| \to \infty\). Then for every \(t\), we have
\begin{align}
    \label{pro1-convolution}
    \int_\mbr \zeta(\sigma+\mmi(t+u))K(u)\mmd u = \sum_{n=1}^\infty \frac{\widehat{K}(\log n)}{n^{\sigma+\mit}}+E(\sigma,t),
\end{align}
where \(E(\sigma,t) = 2\pi K(-t-\mmi(1-\sigma))\).
\end{lem}
\par
The following Lemma shows that the resonator \(R(t)\) constructed in Section \ref{secConstucting} yields an effective lower bound for the contribution from Dirichlet series with non-negative coefficients.

\begin{lem}
    \label{pro2}
    For all \(n \in \mbn_+\), let \(a_n \ge 0\). Assume that \(G(t)=\sum_{n=1}^\infty a_n n^{-1/2-\mit}\) is absolutely convergent. Let \(\varepsilon\) be a small positive number. Then for sufficiently large \(T\), we have
    \begin{align*}
        \int_\mbr\re G(t)& |R(t)|^2\ph\Bgl\frac{t}{T} \Bgr \mmd t  \nonumber \\
        &\ge T \bgl\min_{n \le T^\varepsilon} a_n\bgr \exp \bggl \bgl \eta \sqrt{\kappa} +o(1) \bgr \sqrt{\frac{\log T \lt T}{\lo T}}\bggr \sum_{n \in \mcm}f(n)^2.
    \end{align*}
\end{lem}

\begin{proof}
    Taking the real part of both sides of \cite[Lemma 6, Eq. (23)]{bondarenko2018argument} suffices; we refer to \cite[p. 1699]{bondarenko2017large} for further details.
\end{proof}


\begin{thebibliography}{99}
\bibitem[]{bondarenko2017large} A. Bondarenko; K. Seip. Large greatest common divisor sums and extreme values of the {R}iemann zeta function. Duke Math. J. 166 (2017) 1685--1701
\bibitem[]{bondarenko2018argument} Bondarenko, A; Seip, K. Extreme values of the {R}iemann zeta function and its argument. Math. Ann. 372 (2018) 999--1015
\end{thebibliography}
\end{document}
